\subsection{\rqC}

\subsubsection{Overall Impact of Test Refactorings on Test Smells}

Our analysis reveals that test refactorings have a limited but significant impact on test smells. The descriptive analysis shows that the majority of test smell instances (85.6\% at file level and 91.9\% at refactoring level) remain unchanged after refactoring, indicating that test refactorings typically target specific aspects of test code rather than comprehensively addressing all test quality issues.


\input{200_Tab/RQ3/Summary_TestSmellChanges_AfterRefactoring}

To better understand which test smells are more likely to improve through refactoring, Figures~\ref{fig:test-smell-improvement-rate-file} and~\ref{fig:test-smell-improvement-rate-refactoring} show the improvement rates for each test smell at the \textbf{file level} and \textbf{refactoring level}, respectively. These visualizations help highlight which smells tend to benefit from refactorings and which are more resistant to change.

\begin{figure}[htbp]
\centering
\includegraphics[width=0.85\linewidth]{100_Fig/RQ3/change_patterns_file_better.pdf}
\caption{Improvement rate of each test smell after refactoring (file level)}
\label{fig:test-smell-improvement-rate-file}
\end{figure}

\begin{figure}[htbp]
\centering
\includegraphics[width=0.85\linewidth]{100_Fig/RQ3/change_patterns_method_better.pdf}
\caption{Improvement rate of each test smell after refactoring (refactoring level)}
\label{fig:test-smell-improvement-rate-refactoring}
\end{figure}


\subsubsection{Statistically Significant Improvements}

When focusing on test smell instances that changed after refactoring, we found several statistically significant improvements. Table~\ref{tab:significant-improvements} presents the results of our conditional statistical analysis.

\input{200_Tab/RQ3/sig_testsmell_improvement}


The most notable finding is the substantial improvement in \textit{Duplicate Assert} test smells, which showed large effect sizes at both file level (d = -1.947) and refactoring level (d = -1.212). This suggests that test refactorings are particularly effective at reducing redundant assertions in test code.

\subsubsection{Refactoring Type Effectiveness}

To better understand the relationship between test refactorings and test smells, we analyzed which specific refactoring types are most effective in reducing particular smells. Table~\ref{tab:refactoring-effectiveness} summarizes the refactoring–smell pairs that exhibited the highest improvement rates, highlighting the most effective refactoring types for each test smell.

\begin{table}[htbp]
\centering
\caption{Most Effective Refactoring Types for Test Smell Reduction}
\label{tab:refactoring-effectiveness}
\begin{tabular}{llrrr}
\toprule
\textbf{Test Smell} & \textbf{Best Refactoring Type} & \textbf{Improvement Rate} & \textbf{Effect Size} & \textbf{Significant} \\
\midrule
Duplicate Assert & Split Test Method & 67.6\% & -1.947 & Yes \\
Magic Number Test & Grouping Tests & 69.2\% & -1.123 & Yes \\
Assertion Roulette & Grouping Tests & 46.2\% & -0.234 & No \\
\bottomrule
\end{tabular}
\end{table}

\paragraph{Key Findings:}
\begin{itemize}
    \item \textbf{Split Test Method} is highly effective for reducing \textit{Duplicate Assert} test smells, with a 67.6\% improvement rate and large effect size.
    \item \textbf{Grouping Tests} shows effectiveness across multiple test smells, particularly for \textit{Magic Number Test} (69.2\% improvement rate).
    \item \textbf{Parameterized Test} refactoring demonstrates moderate effectiveness for various test smells.
\end{itemize}

\subsubsection{Limitations and Implications}

Our analysis reveals several important limitations and implications:

\begin{itemize}
    \item \textbf{Limited Scope}: The high proportion of unchanged test smells (85-92\%) suggests that test refactorings address specific, targeted issues rather than comprehensive test quality improvements.
    \item \textbf{Context-Dependent Effectiveness}: The effectiveness of refactorings varies significantly depending on the specific test smell and refactoring type combination.
    \item \textbf{Statistical Power}: Many refactoring-scent combinations have insufficient sample sizes for reliable statistical analysis, indicating the need for larger-scale studies.
\end{itemize}

\paragraph{Answer to RQ3:} Test refactorings do help reduce certain test smells, particularly \textit{Duplicate Assert} and \textit{Magic Number Test}. However, their impact is limited and context-dependent, with most test smell instances remaining unchanged. The effectiveness varies significantly by refactoring type, with \textit{Split Test Method} and \textit{Grouping Tests} showing the most consistent positive effects. 